\section{Local Physics-Based Estimation and Diagnostics}
\label{sec:estimation}

\subsection{Nominal Estimator Hierarchy}

All estimators use the same equilibrium $\bm y_0$, hidden-voltage center $\bm z_{\mathcal U,0}$, sampled dynamics $\bm A_d$, PMU map $\bm C$, and hidden-output map $\bm L_{\mathcal U}$. This shared model isolates the value of measurement updating and temporal propagation.

The equilibrium baseline B0 returns
\begin{equation}
 \widehat{\bm z}^{\mathrm{B0}}_{\mathcal U,k}=\bm z_{\mathcal U,0}. \label{eq:b0}
\end{equation}
It measures how much of the score can be obtained by remaining at the nominal power-flow solution.

The snapshot estimator B1 computes the minimum-variance linear update independently at every frame,
\begin{align}
 \bm K_{\mathrm{WLS}}&=\bm P_0\bm C^{\mathsf T}
 \left(\bm C\bm P_0\bm C^{\mathsf T}+\bm R\right)^{-1}, \label{eq:wls-gain}\\
 \widehat{\delta\bm x}^{\mathrm{B1}}_k&=\bm K_{\mathrm{WLS}}(\bm y_k-\bm y_0), \qquad
 \widehat{\bm z}^{\mathrm{B1}}_{\mathcal U,k}=\bm z_{\mathcal U,0}+\bm L_{\mathcal U}\widehat{\delta\bm x}^{\mathrm{B1}}_k. \label{eq:wls-update}
\end{align}
B1 uses network-constrained cross-covariance but no state carried from frame $k-1$.

The causal Kalman estimator B2 applies the classical linear-Gaussian recursion \cite{kalman1960}:
\begin{align}
 \widehat{\bm x}_{k|k-1}&=\bm A_d\widehat{\bm x}_{k-1|k-1}, \label{eq:kf-predict-state}\\
 \bm P_{k|k-1}&=\bm A_d\bm P_{k-1|k-1}\bm A_d^{\mathsf T}+\bm Q, \label{eq:kf-predict-cov}\\
 \bm\nu_k&=(\bm y_k-\bm y_0)-\bm C\widehat{\bm x}_{k|k-1}, \label{eq:innovation}\\
 \bm S_k&=\bm C\bm P_{k|k-1}\bm C^{\mathsf T}+\bm R, \label{eq:innovation-cov}\\
 \bm K_k&=\bm P_{k|k-1}\bm C^{\mathsf T}\bm S_k^{-1}, \label{eq:kf-gain}\\
 \widehat{\bm x}_{k|k}&=\widehat{\bm x}_{k|k-1}+\bm K_k\bm\nu_k. \label{eq:kf-update}
\end{align}
The covariance update uses the Joseph form. The hidden estimate is $\bm z_{\mathcal U,0}+\bm L_{\mathcal U}\widehat{\bm x}_{k|k}$. B2 is causal because only $\bm y_{0:k}$ enters the result at $k$.

B3 is an exact fixed-lag Gaussian smoother, following the linear Gaussian backward-recursion principle in \cite{rauch1965smoothing}, evaluated at lags $\ell\in\{0,1,3,5,10,15,30,60\}$ frames. It uses measurements through $k+\ell$ to refine the estimate at $k$ and is therefore retrospective by $\ell/30$ seconds. Dense and cached implementations agree numerically on representative lags. B3 tests whether temporal information beyond the causal update improves the hidden-voltage target; it is not used to support a real-time claim.

\subsection{Marginal Calibration and Colored Discrepancy}

For the hidden two-component voltage error at bus $i$, B2 provides
\begin{equation}
 \bm\Sigma_{i,k}=\bm L_i\bm P_{k|k}\bm L_i^{\mathsf T}. \label{eq:hidden-covariance}
\end{equation}
The original covariance is U0. A two-block temperature model U2 is fitted on an independent calibration split:
\begin{equation}
 \bm\Sigma^{\mathrm{U2}}_{i,k}=\bm T\bm\Sigma_{i,k}\bm T^{\mathsf T},\qquad
 \bm T=\operatorname{diag}(\sqrt{\tau_{\Re}},\sqrt{\tau_{\Im}}). \label{eq:temperature}
\end{equation}
The temperatures minimize Gaussian negative log likelihood (NLL) on calibration data and are frozen before TEST. Coverage, NLL, normalized innovation squared (NIS), and temporal innovation correlation are reported separately; nominal marginal coverage alone is not called complete calibration.

The residual study considers a latent Gauss--Markov process
\begin{equation}
 \bm c_{k+1}=\bm F_c\bm c_k+\bm\eta_k. \label{eq:discrepancy-state}
\end{equation}
In D2-M, $\bm c_k$ enters the PMU measurement residual through a calibration-fitted low-rank basis. In D2-P, it enters a sparse physical process dictionary. Basis, rank, $\bm F_c$, and latent covariance are selected using calibration records only. The experiment asks whether the observed low-rank structure improves reconstruction, hidden uncertainty, or innovation whiteness. It does not assume that a lower innovation autocorrelation is sufficient.

\subsection{Model-Adequacy Statistic and Recentring Ablations}

Let $\bm\mu_{\nu}$ and $\bm\Sigma_{\nu}$ be the mean and shrunk covariance of B2 innovations on the nominal calibration set. The per-frame adequacy statistic is the five-frame trailing mean of
\begin{equation}
 a_k=(\bm\nu_k-\bm\mu_{\nu})^{\mathsf T}\bm\Sigma_{\nu}^{-1}(\bm\nu_k-\bm\mu_{\nu}). \label{eq:adequacy-score}
\end{equation}
Mismatch labels and hidden truth do not enter \eqref{eq:adequacy-score}; hidden truth is used only after inference to label frames whose hidden TVE exceeds the frozen evaluation threshold. AUROC and AUPRC therefore measure whether observed innovations anticipate unacceptable hidden error, not whether the detector identifies a mismatch family.

Two evaluation-only or adaptive variants diagnose the failure mechanism. R1 replaces the nominal equilibrium with the true plant equilibrium and is explicitly oracle. R2-A estimates an eight-dimensional grouped observed-bus voltage offset using a 60-frame trailing window, nominal matrices, ridge regularization, and only PMU measurements. It is a preregistered linearized static-MAP proxy, not a nonlinear AC power-flow or KKT solve. Its center/state cross-covariance is neglected. Comparing frozen B2, R1, and R2-A tests whether a simple causal offset correction captures the oracle recentering benefit.
